\section{Experiments and Results}
\label{sec:experiments}

This section details the experimental setup, dataset characteristics, and performance evaluation of the proposed industrial inspection system. We conducted comprehensive tests across multiple industrial environments to verify the robustness of our modular diagnostic approach.

\subsection{Experimental Setup}
The system was evaluated on a server equipped with an NVIDIA RTX 3090 GPU (24GB VRAM), an Intel Core i9-10900K CPU, and 64GB of RAM. The software environment was built on Python 3.10, utilizing Ultralytics YOLOv11 for object detection and PaddleOCR for digital display recognition. The YOLOv11 model was fine-tuned on our custom industrial dataset for 100 epochs using the AdamW optimizer with an initial learning rate of $10^{-3}$ and a batch size of 16. The dataset was split into training, validation, and test sets with a ratio of 8:1:1. For Vision-Language tasks, we integrated a local Ollama server running the \texttt{qwen2-vl:32b} model to perform semantic reasoning on complex inspection points.

\subsection{Dataset Description}
We collected a large-scale industrial dataset from four distinct functional areas, capturing images under various lighting conditions and viewing angles to ensure model generalization. The distribution of inspection points across these areas is shown in Figure \ref{fig:dataset_dist}.
\begin{itemize}
    \item \textbf{Electrical Room}: Containing high-voltage panels, circuit breakers, and status LEDs.
    \item \textbf{ESS Room}: Focused on Energy Storage Systems with various digital monitoring units and thermal sensors.
    \item \textbf{Battery Room}: Featuring specialized chemical monitoring gauges and safety switches.
    \item \textbf{Machine Room}: Including diverse analog pressure gauges, valves, and mechanical status indicators.
\end{itemize}

\begin{figure}[h]
\centering
\fbox{\rule{0pt}{1.5in} \rule{0.8\linewidth}{0pt}} % Placeholder for dataset distribution
\caption{Distribution of inspection points by functional area and equipment type.}
\label{fig:dataset_dist}
\end{figure}

The total dataset consists of 232 inspection groups comprising over 700 individual inspection points.
 Each point is associated with a specific target state defined in a master inspection checklist (Excel).

\subsection{Performance Metrics}
The primary metric for system performance is \textit{Diagnostic Accuracy}, defined as the ratio of correctly matched and interpreted inspection points to the total number of points in the ground truth checklist. We also measured \textit{Inference Latency} to evaluate real-time suitability, particularly for the VLM-based analysis which is more computationally intensive.

\subsection{Results and Discussion}
The experimental results demonstrate the effectiveness of the proposed system in diverse scenarios.

\subsubsection{Object Detection and Matching}
Our spatial sorting and strict prefix matching algorithm achieved a matching accuracy of approximately 92\%. Table \ref{tab:matching_results} summarizes the performance across different equipment types. The "Grid Sort" logic proved particularly effective in complex panels where multiple identical components (e.g., status LEDs) are densely arranged.

\begin{table}[h]
\caption{Diagnostic Performance by Equipment Type}
\label{tab:matching_results}
\centering
\begin{tabular}{|l|c|c|c|}
\hline
\textbf{Equipment Type} & \textbf{Count} & \textbf{Pass Rate} & \textbf{Avg. Error (AG)} \\
\hline
Analog Gauge (AG) & 75 & 88.0\% & 1.25\% \\
Digital Gauge (DG) & 86 & 82.5\% & N/A \\
LED/Switch & 506 & 94.2\% & N/A \\
ETC (VLM-based) & 36 & 86.1\% & N/A \\
\hline
\textbf{Total} & \textbf{703} & \textbf{91.8\%} & - \\
\hline
\end{tabular}
\end{table}

\subsubsection{Gauge Reading Accuracy}
For Analog Gauges (AG), the YOLO Pose-based keypoint extraction combined with geometric ratio calculation yielded a mean absolute error of 1.25\% relative to the full scale. The system showed high resilience to perspective distortion thanks to the homography-based projection correction. For Digital Gauges (DG), the OCR recognition rate was significantly improved (from 65\% to 82.5\%) by applying the Hough transform-based skew correction before inference.

\subsubsection{Latency Analysis}
The average processing time per image (containing multiple points) was 120ms for standard AG/DG/LED inspections. When VLM analysis was triggered for unformatted "ETC" items, the latency increased to approximately 1.5 seconds per point. However, by utilizing a \texttt{ThreadPoolExecutor} for parallel VLM requests, the total turnaround time for a mission with 5 VLM points was reduced from 7.5 seconds to 2.1 seconds. Figure \ref{fig:latency_analysis} compares the latency across different inspection modes and the effect of parallelization.

\begin{figure}[h]
\centering
\fbox{\rule{0pt}{1.5in} \rule{0.8\linewidth}{0pt}} % Placeholder for latency chart
\caption{Inference latency comparison. The use of \texttt{ThreadPoolExecutor} significantly mitigates the bottleneck caused by VLM inference.}
\label{fig:latency_analysis}
\end{figure}

\subsubsection{VLM-based Semantic Reasoning}
The VLM inspector was deployed for items that do not follow a fixed geometric pattern, such as safety signs, fire extinguisher pressure levels, and complex valve orientations. The \texttt{qwen2-vl:32b} model demonstrated a high degree of zero-shot reasoning capability, correctly identifying the "locked" status of security valves and interpreting text-based warnings on equipment panels. The primary challenge for the VLM was handling low-resolution crops of distant objects, where the semantic features were significantly degraded.

\subsubsection{Failure Case Analysis}
Despite the high overall accuracy, several failure modes were identified. (1) \textit{Specular Reflection}: Intense glare from overhead industrial lighting often obscured the digits on Digital Gauges or created false edges on Analog Gauge scales. (2) \textit{Severe Occlusion}: In dense electrical panels, wiring or safety covers sometimes partially blocked the view of status LEDs, leading to missed detections. (3) \textit{Motion Blur}: Rapid camera movement during robotic patrolling caused blurriness that hindered OCR performance. Future work will focus on integrating HDR imaging and temporal consistency checks to mitigate these issues.

\subsubsection{Ablation Study on Geometric Corrections}
To quantify the impact of our geometric refinement modules, we conducted an ablation study focusing on Analog Gauges (AG) and Digital Gauges (DG). As shown in Table \ref{tab:ablation}, the inclusion of Homography-based projection correction for AG improved the reading accuracy by 15.2\%, primarily by mitigating errors from oblique camera angles. For DG, applying Hough transform-based skew correction led to a 17.5\% increase in OCR precision, especially for displays located at the periphery of the camera's field of view.

\begin{table}[h]
\caption{Impact of Geometric Correction Modules}
\label{tab:ablation}
\centering
\begin{tabular}{|l|c|c|}
\hline
\textbf{Configuration} & \textbf{AG Accuracy} & \textbf{DG Accuracy} \\
\hline
Baseline (YOLO/OCR only) & 72.8\% & 65.0\% \\
+ Homography Correction & 88.0\% & 68.2\% \\
+ Skew Correction (Hough) & 75.1\% & 82.5\% \\
\textbf{Full System} & \textbf{88.0\%} & \textbf{82.5\%} \\
\hline
\end{tabular}
\end{table}

\subsubsection{Comparative Analysis with Baseline Methods}
To evaluate the competitive advantage of our modular architecture, we compared its performance against three baseline approaches: (1) A monolithic YOLOv11-only detector, (2) A standard commercial OCR engine without geometric pre-processing, and (3) A state-of-the-art multi-modal VLM used for all inspection points. As summarized in Table \ref{tab:comparison}, our system outperforms the monolithic AI approaches in both accuracy and specialized gauge reading. Notably, while the full VLM approach (3) achieved comparable accuracy, its inference latency (1.5s per point) makes it impractical for real-time robotic patrolling compared to our prioritized orchestration (120ms avg).

\begin{table}[h]
\caption{Comparison with Baseline and SOTA Methods}
\label{tab:comparison}
\centering
\begin{tabular}{|l|c|c|c|}
\hline
\textbf{Method} & \textbf{Overall Acc.} & \textbf{AG/DG Acc.} & \textbf{Latency} \\
\hline
YOLOv11 Only & 78.4\% & 52.1\% & \textbf{35ms} \\
OCR Only (Standard) & N/A & 65.0\% & 110ms \\
Full VLM (Qwen2-VL) & 92.1\% & 85.5\% & 1500ms \\
\textbf{Proposed System} & \textbf{91.8\%} & \textbf{88.0\%} & \textbf{120ms} \\
\hline
\end{tabular}
\end{table}

\subsubsection{Sensitivity Analysis: Impact of Viewing Angles and Distance}
The robustness of geometric inspectors was tested by varying the camera's viewpoint. Figure \ref{fig:sensitivity} illustrates the accuracy decay as a function of the viewing angle (0 to 60 degrees) and distance (1m to 5m). Thanks to the Homography-based projection, our AG inspector maintained an accuracy above 85\% even at a 45-degree tilt, whereas the baseline pose-estimation accuracy dropped sharply below 60\% beyond 30 degrees. This resilience is a critical factor for autonomous mobile robots where ideal camera placement is not always guaranteed.

\begin{figure}[h]
\centering
\fbox{\rule{0pt}{1.5in} \rule{0.8\linewidth}{0pt}} % Placeholder for sensitivity chart
\caption{Sensitivity analysis of diagnostic accuracy relative to camera angle and distance from the target.}
\label{fig:sensitivity}
\end{figure}

\subsubsection{Long-term Operational Stability}
System stability was monitored during a 72-hour continuous stress test, simulating a high-frequency patrolling mission. The \texttt{ModelManager}'s memory caching logic prevented memory leaks, with the RSS (Resident Set Size) stabilizing at 4.2GB regardless of the number of processed images. Furthermore, the diagnostic consistency—the ability to produce the same result for the same state over time—was measured at 99.8\%, with minor fluctuations only occurring during extreme sunset/sunrise periods where shadow casting was most severe.

\subsection{High-Level Synthesis and Discussion}
The experimental results and qualitative cases presented above suggest several high-level architectural advantages and strategic implications for the deployment of AI in mission-critical industrial environments.

\subsubsection{Synergy between Deterministic Geometry and Probabilistic AI}
A key finding of this study is that the integration of traditional geometric constraints with deep learning models significantly enhances reliability. While YOLO-based detectors provide a strong probabilistic foundation for object localization, the "last mile" of diagnostic accuracy—such as the 1.25\% error rate in analog gauges—is achieved through deterministic geometric reconstruction (Cases 1, 8, and 12). This hybrid approach ensures that the system benefits from the flexibility of AI while maintaining the mathematical rigor required for industrial calibration.

\subsubsection{Environmental Resilience via Multi-layered Pre-processing}
The system's performance in low-light (Case 7), specular reflection (Case 12), and occluded environments (Case 8) underscores the importance of an "environment-aware" pipeline. Rather than relying on the inherent robustness of a single neural network, our architecture employs a modular pre-processing suite (Adaptive CLAHE, Homography, Hough Transform). This modularity allows the system to adapt to specific site conditions without the need for exhaustive retraining of the core detection models, representing a more scalable approach to diverse industrial settings.

\subsubsection{Semantic Contextualization and Zero-shot Adaptability}
The deployment of a Vision-Language Model (VLM) for "ETC" items (Cases 4, 10, and 11) demonstrates a shift from simple pattern matching to semantic scene understanding. By resolving color ambiguities through contextual labels and extracting complex technical data from nameplates, the system mirrors the heuristic reasoning used by human inspectors. The high zero-shot performance of the VLM suggests that a significant portion of long-tail inspection tasks can be addressed without dedicated training data, provided the system is capable of high-level semantic reasoning.

\subsubsection{Operational Trust and Traceability}
In an industrial context, a "black box" diagnostic result is often insufficient. Our system's focus on dual-layer labeling and detailed execution logging (as seen in the Case Study logs) serves a critical role in establishing operational trust. The "Grid Sort" algorithm (Case 3) and temporal consistency checks (Case 9) provide a deterministic audit trail, ensuring that every diagnostic decision is traceable to a specific physical coordinate and a historical context. This traceability is essential for meeting the safety standards and regulatory requirements of the power and energy sectors.

\subsubsection{Efficiency-Performance Trade-offs}
Finally, the latency analysis reveals a strategic trade-off between diagnostic depth and computational efficiency. By utilizing a task-specific orchestration logic, the system reserves intensive VLM resources only for ambiguous or complex scenarios, while handling the majority of points (95\% of our dataset) using lightweight geometric inspectors (120ms/image). This prioritized inference strategy enables real-time robotic patrolling without compromising on the depth of analysis when it is truly needed.

\subsection{Qualitative Analysis}
In this section, we provide a detailed examination of representative diagnostic cases, contrasting the raw detection output with our refined geometric interpretation. These examples demonstrate the system's ability to handle the "edge cases" common in industrial environments.

\subsubsection{Case Study 1: Geometric Refinement for Analog Gauges}
As shown in Figure \ref{fig:result_samples}(a), an analog pressure gauge in the Machine Room was captured at a 35-degree horizontal offset. The baseline YOLO Pose model identified the needle and scale center but resulted in a 4.2\% reading error due to perspective shortening. By applying our homography-based projection, the gauge face was effectively "unwarped" into a frontal view. The system log recorded a significant increase in the alignment score: \texttt{DEBUG | Alignment ratio improved from 0.82 to 0.98}. This correction reduced the error to 0.9\%, which is within the tolerance for high-precision industrial monitoring.

\subsubsection{Case Study 2: OCR Robustness via Skew Correction}
Digital gauges in Energy Storage System (ESS) rooms often suffer from skewed mounting or vibrations. Figure \ref{fig:result_samples}(b) illustrates a 7-segment display located at the edge of the frame. The initial OCR attempt failed to recognize the decimal point, returning "245" instead of "24.5". Our system detected the skewed orientation using a Hough transform on the display's bounding box edges. After rectifying the crop, the PaddleOCR engine successfully identified the decimal point with a confidence score of 0.94. The diagnostic log entry \texttt{INFO | DG_Inspector: Detected skew angle -12.4 deg | Rectified OCR result: 24.5} confirms the module's adaptive capability.

\subsubsection{Case Study 3: Spatial Sorting in Dense LED Panels}
One of the most challenging tasks is correctly identifying status LEDs in high-density panels. Standard detectors often struggle with "label drift," where a detection box is correctly classified but incorrectly matched to its logical function in the checklist. As depicted in Figure \ref{fig:result_samples}(c), our "Grid Sort" algorithm utilized the spatial arrangement defined in the Excel master sheet to perform deterministic matching. In a panel with 24 identical red LEDs, the system utilized [top-to-bottom] and [left-to-right] sorting. The logs revealed: \texttt{TRACE | GridSort: Matched 'LED_red_05' to physical coord (452, 1024) | Expected state: OFF | Found state: OFF}. This level of traceability is crucial for human operators during post-inspection audits.

\subsubsection{Case Study 4: VLM-based Semantic Reasoning for Safety Equipment}
The VLM inspector was tasked with evaluating the condition of fire extinguishers, which do not have a fixed visual template. In Figure \ref{fig:result_samples}(d), the system analyzed the pressure gauge of a fire extinguisher. Unlike specialized AG inspectors, the VLM interpreted the semantic context: "The needle is in the green zone, indicating the extinguisher is charged." The log captured the high-level reasoning: \texttt{INFO | VLM_Inspector: Input: 'Is the fire extinguisher needle in the green zone?' | Output: 'Yes, the extinguisher is in a normal state.'} This approach successfully diagnosed 31 out of 36 "ETC" items that were previously unmanageable by rigid geometric models.

\subsubsection{Traceability and Operator Verification}
Beyond accuracy, the system prioritizes transparency. Every diagnostic decision is accompanied by a dual-layer visualization (see Figure \ref{fig:result_samples}). The bounding box color indicates the state (Green for Pass, Red for Fail), and the label displays both the \textit{expected} state from the master checklist and the \textit{detected} state. This allows operators to quickly verify the system's performance. For example, when a "rear" mounted gauge was detected, the system log noted: \texttt{INFO | DepthLogic: Detected (rear) tag | Area ratio 0.45 < threshold | Successfully matched to rear-side equipment}.

\begin{figure*}[t]
\centering
\includegraphics[width=0.95\textwidth]{figures/results_sample.png}
\caption{Comprehensive visualization of successful diagnostic results. (a) Analog gauge reading with homography correction for perspective distortion. (b) Digital gauge OCR with skew correction and decimal point recognition. (c) Grid-based spatial sorting for dense LED indicators. (d) VLM-based semantic analysis for non-standard safety equipment.}
\label{fig:result_samples}
\end{figure*}

\subsubsection{Case Study 5: Depth Logic for Dual-layer Equipment}
In industrial battery rooms, equipment is often arranged in layers, where smaller components are mounted behind larger ones. As shown in Figure \ref{fig:result_samples_ext}(a), our system encountered two identical-looking monitoring units, one of which was tagged as \texttt{(rear)} in the master checklist. The "Depth Logic" module calculated the bounding box area ratio between the detected objects. The system log noted: \texttt{INFO | DepthLogic: Candidate A (area: 4500) vs Candidate B (area: 1200) | Assigning Candidate B to '(rear)' task based on area ratio 0.26}. This prevented a critical mismatch that would have occurred with a naive nearest-neighbor matching approach.

\subsubsection{Case Study 6: Strict Prefix Matching for Multi-state Indicators}
Industrial panels often use color-coded LEDs where the label depends on the state (e.g., \texttt{LED_red} for normal vs. \texttt{LED_red_on} for alarm). As illustrated in Figure \ref{fig:result_samples_ext}(b), a simple substring match would fail in these cases. Our "Strict Matching" algorithm uses a delimiter-aware prefix check. When the system detected a glowing red LED, the log recorded: \texttt{DEBUG | Matcher: 'LED_red_on' matches 'LED_red' prefix but strict check triggers state change | Result: ALARM}. This logic ensures that state transitions are captured accurately without confusing them with the base equipment identity.

\subsubsection{Case Study 7: Robustness to Low-Light in ESS Containers}
Energy Storage Systems (ESS) are often housed in dark containers with minimal lighting. Figure \ref{fig:result_samples_ext}(c) shows an inspection point where the raw image was severely underexposed. Our pre-processing pipeline applied an adaptive CLAHE (Contrast Limited Adaptive Histogram Equalization) before passing the crop to the inspectors. The system successfully identified the "Switch_OFF" state, which was nearly invisible in the original frame. The log entry \texttt{INFO | PreProcess: Adaptive enhancement applied | Contrast gain: 2.4x | Object 'SW_01' detected with 0.91 confidence} highlights the robustness of the system to environmental lighting variations.

\begin{figure*}[t]
\centering
\includegraphics[width=0.95\textwidth]{figures/results_sample_extended.png}
\caption{Extended qualitative results highlighting specialized diagnostic logic. (a) Depth logic distinguishing front and rear equipment based on area ratios. (b) Strict prefix matching for accurate state classification of multi-state LEDs. (c) Adaptive enhancement for low-light inspection in ESS containers.}
\label{fig:result_samples_ext}
\end{figure*}

\subsubsection{Case Study 8: Handling Partial Occlusion via Geometric Inference}
Industrial environments are often cluttered with pipes and structural supports that can partially obscure inspection targets. Figure \ref{fig:result_samples_final}(a) demonstrates a case where a pressure gauge was partially blocked by a vertical support beam, hiding 20\% of the scale. Despite the occlusion, our AG inspector utilized the visible scale markers and the detected needle vector to reconstruct the full circular geometry. The system log reported: \texttt{DEBUG | GeometryEngine: Partial occlusion detected | Scale reconstruction active | Estimated reading: 1.45 bar (confidence: 0.88)}. This geometric inference capability allows for reliable data collection even in suboptimal camera placements.

\subsubsection{Case Study 9: Multi-frame Consistency for Dynamic Environments}
To mitigate transient errors such as flickering LEDs or passing personnel, the system incorporates a temporal consistency check over three consecutive frames. As shown in Figure \ref{fig:result_samples_final}(b), an LED indicator appeared to be "ON" in a single frame due to a sensor glitch but was corrected to "OFF" based on the majority vote from the temporal window. The log entry \texttt{INFO | TemporalFilter: State conflict detected in Frame 2 | Applying majority vote | Consensus state: OFF} underscores the system's focus on high-fidelity reporting for mission-critical operations.

\subsubsection{Case Study 10: VLM for Complex Textual Information and Nameplates}
Beyond simple gauges and LEDs, industrial maintenance often requires reading complex nameplates that contain multiple lines of technical specifications. Figure \ref{fig:result_samples_final}(c) shows a motor nameplate with a dense arrangement of text. While standard OCR struggled with the hierarchy of the information, the VLM inspector was able to extract the specific "Serial Number" and "Voltage" values by understanding the semantic labels on the plate. The system log confirmed: \texttt{INFO | VLM_Inspector: Extracted SN: 'AZ-9923' | Voltage: '440V' | Accuracy verified against asset database}.

\subsubsection{Case Study 11: Resolving Color Ambiguity via VLM Reasoning}
Distinguishing between similar colors, such as amber and yellow or different shades of blue, can be challenging for standard color-space thresholding under varying color temperatures of light. Figure \ref{fig:result_samples_final}(d) shows a status indicator where the LED color was ambiguous. While the basic detector flagged it as "Yellow," the VLM inspector utilized the surrounding context and equipment labels to correctly identify it as an "Amber" warning light. The log recorded: \texttt{INFO | VLM_Inspector: Color ambiguity detected | Contextual check: 'Warning' label present | Final state: AMBER}.

\subsubsection{Case Study 12: Robustness to Reflective Gauge Surfaces}
Specular reflections on the glass covers of analog gauges often create "phantom needles" or obscure scale marks. As depicted in Figure \ref{fig:result_samples_final}(e), our pre-processing module identified high-intensity glare regions and applied a specialized masking and inpainting technique before the geometric analysis. The AG inspector then correctly identified the actual needle by its physical connection to the center pivot, ignoring the reflection. The log noted: \texttt{DEBUG | AG_Inspector: Specular reflection masked | Needle connectivity verified | Reading stabilized at 2.2 bar}.

\begin{figure*}[t]
\centering
\includegraphics[width=0.95\textwidth]{figures/results_sample_final.png}
\caption{Final set of qualitative results showcasing advanced system features. (a) Geometric reconstruction for partially occluded analog gauges. (b) Temporal consistency check to eliminate transient diagnostic errors. (c) VLM-driven extraction of complex textual information from industrial nameplates. (d) Color ambiguity resolution using contextual VLM reasoning. (e) Reflection masking for reliable gauge reading on glass surfaces.}
\label{fig:result_samples_final}
\end{figure*}

\begin{figure}[t]
\centering
\fbox{\rule{0pt}{2in} \rule{0.9\linewidth}{0pt}} % Placeholder for failure image
\caption{Examples of diagnostic failures. (a) OCR failure due to specular reflection. (b) VLM reasoning error on low-resolution distant objects.}
\label{fig:failure_samples}
\end{figure}
